\section{Verification and formalization plan}

\subsection{What was checked}

The source statements were read at the pinned revision and compared
directly with the hypotheses used in the proofs. In particular, the
near-isometry condition quantifies over all source subspaces; $s>0$ is
explicit; the target projection commutes with the group action; the two
raw errors apply to the same scaled Poincare map; source coordinates and
nonzero normalizing factors preserve image dimensions; and the
ramification loss is charged once per selected translate.

Three exact checkers accompany the arguments. The two
spanning checkers use rational arithmetic for the integer rank schedule,
498 cyclic models, equality examples, and constant checks. The $S_5$ checker generates the
specified subgroups, verifies the character relations, sums the
cyclotomic values exactly, and checks the induced-character identities
on every conjugacy type. Its conclusions are finite algebraic facts;
the passage from those identities to holomorphy is the written
Artin--Hecke argument.

The PDF was compiled from the enclosed TeX and inspected after rendering.
The full upstream Lean project was not built, and no new theorem here is
labelled Lean-verified. A successful script run supports its stated
finite checks; it is not a substitute for the general proofs.
Commands and dependencies are in \code{README.md};
\code{checks/run_checks.py} replays all finite checks offline.
The source excerpts are not a complete upstream Lean checkout.

\subsection{A concrete Lean implementation route}

The most direct formal addition is a separate-error version of
\code{New.isSpread_of_second_moment}. Reuse
\code{sq_sum_hilbertSchmidtNorm_orthoProj_comp_le}, retain the square root
instead of the arithmetic--geometric mean relaxation, and carry the
near-isometry parameter $u$ independently of the moment excess $v$.
The result should quantify over each target submodule and conclude the
preimage-rank inequality in Lemma~\ref{lem:sqrt}.

Next, reuse \code{finrank_sup_range_add_finrank_comap} to formalize the
greedy recurrence, with the stopped case handled explicitly when the
codimension is below $d\delta_*$. The max in $q=\max(0,1-h)$ avoids an
unjustified positivity assumption on $1-h$. A separate lemma can implement
the source-restriction loss, reusing the existing finite-supremum rank
inequality. The final specialization should invoke the two raw source
declarations before their enlargement in \code{isSpread_goodPoincare}.

The following map identifies the source declarations needed for review.
The full relative paths and hashes are included in the evidence manifest.

\begin{center}\small
\begin{tabularx}{\linewidth}{@{}p{0.26\linewidth}X@{}}
\toprule
Role & Declaration or file\\
\midrule
Near-isometry definition & \code{New.IsNearIsometry}, in \code{FewTranslates.lean}\\
Raw near-isometry & \code{Data.isNearIsometry_e}\textsubscript{0}\code{_comp_scaledPoincare}\\
Raw second moment & \code{Data.targetDim_mul_sum_hilbertSchmidtNorm_sq_le}\\
Coordinate transport & \code{Data.comap_scaledPoincare_eq_comap_goodPoincare}\\
Normalized images & \code{Data.range_normalizedPoincare}\\
Current target theorem & \code{Data.enough_forms}, in \code{ShortAndEnough.lean}\\
Use of the target theorem & \code{New.star_inequality}, in \code{MainInequality.lean}\\
\bottomrule
\end{tabularx}
\end{center}

The $S_5$ finite identity has a separate formalization route through
permutations, subgroup-generated closures, cyclotomic linear characters,
and the induction formula. The Artin--Hecke continuation theorem is a
much larger library dependency. Keeping these two stages distinct makes
the exact scope of a future formal result clear.

